\section{行列式与 log-determinant 的微分}
\label{sec:10-Matrix-Calculus/03-Trace-Calculus/03}

你在优化一个协方差矩阵 \(\Sigma\)，目标是最小化某种损失。梯度下降需要 \(\nabla_\Sigma\log\det\Sigma\)。你盯着 \(\det\Sigma\) 看了半天——它是 \(\Sigma\) 的所有特征值的乘积。一行行求导？一个 \(d\times d\) 矩阵有 \(d^2\) 个独立参数，展开成多项式再逐项微分，手算要崩溃。更要命的是数值：一个 \(100\times100\) 的正定矩阵，特征值从 0.01 到 50，乘积要么上溢要么下溢。理论推导里行列式到处出现，数值计算里它却是个危险的原料。

这就是为什么你需要的不是行列式的微分，而是 \(\log\det X\) 的微分——对数把乘法变成加法，同时顺手把梯度化成 trace 形式。

\subsection{Jacobi 公式：行列式微分只用一行 trace}

设 \(X\) 可逆。取一个小扰动 \(X\mapsto X+dX\)，Jacobian 公式给出

\[
d\det X = \det X\cdot\trace(X^{-1}dX).
\]

推导思路分两步。首先

\[
\det(X+tH)=\det X\cdot\det(I+tX^{-1}H).
\]

然后展开 \(\det(I+tM)\)。把 \(M\) 的特征值记作 \(\lambda_i\)，则

\[
\det(I+tM)=\prod_i(1+t\lambda_i)=1+t\sum_i\lambda_i+O(t^2)=1+t\trace M+O(t^2).
\]

一阶项就是 trace。因此

\[
d\det X = \det X \cdot \trace(X^{-1}dX).
\]

立即得到 log-determinant 的微分：

\[
d\log\det X = \frac{1}{\det X}d\det X = \trace(X^{-1}dX).
\]

这比直接微分行列式舒服太多了：去掉了前面的 \(\det X\) 因子，只剩一个 trace。从微分形式反推梯度：

\[
\trace(X^{-1}dX) = \trace\big((X^{-\trans})^{\trans}dX\big),
\]

所以

\[
\nabla_X\log\det X = X^{-\trans}.
\]

进而

\[
\nabla_X\det X = \det X \cdot X^{-\trans}.
\]

一元退化检验：\(d=1\) 时行列式就是标量本身，上面给出 \(\frac{d}{dx}\log x=1/x\)，行列式梯度给出 \(x\cdot 1/x=1\)。转置记号在标量情形自动消失，一切吻合。

一个关键计算习惯：\(\log\det X\) 的微分永远只出 trace，而不是矩阵本身。trace 可以用循环置换变形、可以与 Frobenius 内积对接、可以用 Hutchinson 估计对巨大矩阵做随机近似。打开 trace 这扇门，行列式的微分就从代数重体力活变成了迹计算的本职工作。

\subsection{log-det 在正定锥上是凹的}

令 \(X\succ0\)，沿对称方向 \(H\)（保持 \(X+tH\) 始终对称的扰动）考察

\[
\phi(t)=\log\det(X+tH).
\]

一阶导数就是 Jacobi 公式的直接应用：

\[
\phi'(0)=\trace(X^{-1}H).
\]

二阶需要把逆的微分 \(d(X^{-1})=-X^{-1}(dX)X^{-1}\) 用上：

\[
\phi''(0) = -\trace(X^{-1}HX^{-1}H).
\]

把 \(H\) 用 \(X^{1/2}\) 做合同变换：令 \(S=X^{-1/2}HX^{-1/2}\)，则

\[
\phi''(0) = -\trace(S^2).
\]

因为 \(H\) 是对称扰动，\(S\) 也是对称的，所以

\[
\trace(S^2)=\trace(S^{\trans}S)=\norm{S}_F^2\ge0,
\]

进而

\[
\phi''(0) = -\norm{S}_F^2 \le 0.
\]

\(\log\det X\) 的二阶导数恒非正——它在整个正定锥上是凹的。注意这里的扰动方向必须对称，否则 \(X+tH\) 不再是协方差/精度矩阵的合法候选，Hessian 形式也会不同。

凹性带来的推论：\(-\log\det X\) 是凸的，并且当 \(X\) 逼近奇异边界时 \(-\log\det X\to+\infty\)。这意味着 \(-\log\det X\) 是半正定锥内部的障碍函数——它在边界处立起无穷高的墙，任何从内部出发的优化算法只要极小化含 \(-\log\det\) 的目标，就不会跨过奇异边界。许多半正定规划的内点法正是靠 log-det 障碍函数把约束``编码''进目标。

\subsection{实验设计的体积直觉}

借助同一个体积解释，log-det 在实验设计中有一个漂亮的物理直觉。设第 \(i\) 类测量提供方向 \(a_i\)，投入权重 \(w_i\ge0\) 后的信息矩阵为

\[
M(w)=\sum_i w_i a_i a_i^{\trans}.
\]

在预算约束下最大化 \(\log\det M(w)\)，等于同时拉伸所有参数方向的信息，使置信椭球体积最小——在线性 Gaussian 近似中，\(\log\det M(w)^{-1}\) 正比于置信椭球体积的对数。

若所有传感器几乎沿同一方向测量，\(M(w)\) 接近奇异，\(\log\det\) 的极小值直接暴露仍有方向未被识别。但准则只评估局部线性化下的信息：共同偏差、错误的噪声模型或未纳入的物理变量，不会因为行列式变大了就自动消失。行列式最大化的输出不是终极真理，而是一个在给定线性框架里信息覆盖最均匀的资源分配方案。

\subsection{链式法则与复合矩阵函数}

当 \(X\) 本身是参数的函数 \(X=X(\theta)\) 时，链式法则给出：

\[
\frac{d}{d\theta}\log\det X(\theta) = \trace\left(X(\theta)^{-1}\frac{dX}{d\theta}\right).
\]

与上一节对逆的微分一样，计算时不需要显式求 \(X^{-1}\) 再乘矩阵。若需要精确的 trace 值，可用 \(X^{-1}\) 的分解（如 Cholesky）分别求解多个右端再累积对角线贡献。若矩阵巨大、只要无偏估计，可用 Hutchinson 估计：

\[
\trace(M) = \E[z^{\trans}Mz],
\]

其中 \(\E[zz^{\trans}]=I\)。取几个随机向量即可估计 trace，代价是引入采样误差——在深度学习的大矩阵优化中，这种随机 trace 估计常是用得上的工程技巧。

Gaussian 对数似然的典型项：

\[
\frac12\log\det\Sigma + \frac12(x-\mu)^{\trans}\Sigma^{-1}(x-\mu).
\]

前半段惩罚总体体积（协方差不能缩成一个点），后半段衡量 Mahalanobis 距离。两者的导数都可通过 Cholesky 分解和回代求解计算，不需要显式求逆——跟前面逆的微分的结论完全衔接。

\subsection{数值上的 log-determinant 计算}

对对称正定矩阵 \(X=LL^{\trans}\)：

\[
\log\det X = 2\sum_i\log L_{ii}.
\]

用到了 \(\det X=(\det L)^2\) 和三角矩阵行列式等于对角元乘积。这避免了大批特征值相乘再取对数——Cholesky 分解本就是解线性系统和计算二次型所需要的，顺带取对角元的对数求和，零额外代价。

一般可逆矩阵的行列式可能为负，实数里 \(\log\det X\) 未定义。软件通常返回符号和 log-absolute-determinant 的分离表示：

\[
\det X = s\,e^\ell,\qquad s\in\{-1,1\},\quad \ell=\log|\det X|.
\]

不要把 \(\log|\det X|\) 和 \(\log\det X\) 混为一体。前者定义域更宽（只要 \(|X|\neq0\) 即可），但两者在负行列式的几何解释和某些导数推导上不可互换。写代码时，如果你调用的是 `slogdet`，返回的就是 \(\log|\det X|\) 加符号位。

当矩阵奇异时，\(\det X=0\)，\(\log\det\) 趋于 \(-\infty\)，逆的公式不再适用。接近奇异时，梯度 \(X^{-\trans}\) 的元素可能巨大——这不是公式的缺陷，它准确反映了体积对最小特征值的极端敏感性。一个接近退化的变换，稍微动一下参数就能把体积推出几个数量级，梯度大是诚实信号：这里对参数极其敏感，你的优化可能正在一片数值峭壁上行走。
